\begin{solapartat}
\figura[0.6]{sol_fig1.png}

Del gràfic es pot extraure: $T = 1,25\ \mathrm{ms}$ \punts{0,1}, $A = 1,5\ \mathrm{cm}$ \punts{0,1}

\punts{0,1} $f = \dfrac{1}{T} = 800\ \mathrm{s^{-1}}$

\punts{0,1} $v = \lambda f = 160\ \mathrm{m\,s^{-1}}$

\punts{0,2} $y(x, t) = A\cos(\omega t - kx + \phi_0)$

Si s'usa la funció cosinus, $\phi_0 = 0$. Si s'usa la funció sinus, aleshores
\[ \phi_0 = \frac{\pi}{2}\ \mathrm{rad} \qquad \text{o bé} \qquad \phi_0 = -\frac{3\pi}{2}\ \mathrm{rad}. \]

\punts{0,1} $\omega = 2\pi f = 1600\pi\ \mathrm{rad/s}$

\punts{0,1} $k = \dfrac{2\pi}{\lambda} = 10\pi\ \mathrm{rad/m}$

\punts{0,2} $y(x, t) = 1,5\ \mathrm{cm}\,\cos\left(1600\pi\,\dfrac{\mathrm{rad}}{\mathrm{s}}\,t - 10\pi\,\dfrac{\mathrm{rad}}{\mathrm{m}}\,x\right)$

També serien vàlides expressions equivalents (amb sinus, amb $A$ en metres, traient el $\pi$ factor comú, escrivint $\omega = 5027\ \mathrm{rad/s}$ i $k = 31,4\ \mathrm{rad/m}$, etc.
\end{solapartat}

\begin{solapartat}
\punts{0,4} $v_y(x, t) = \dfrac{dy}{dt} = -A\omega\sin(\omega t - kx + \phi_0)$

\punts{0,6} Quan el valor de la fase $(\omega t - kx + \phi_0)$ va prenent tots els valors possibles, el seu sinus varia entre $-1$ i $+1$. El valor màxim de la velocitat és $A\omega$.
\end{solapartat}
